% A LaTeX (non-official) template for ISAE projects reports
% Copyright (C) 2014 Damien Roque
% Version: 0.2
% Author: Damien Roque <damien.roque_AT_isae.fr>

\documentclass[a4paper,12pt]{book}
\usepackage[utf8]{inputenc}
\usepackage[T1]{fontenc}
\usepackage[frenchb]{babel} % If you write in French
%\usepackage[english]{babel} % If you write in English
\usepackage{a4wide}
\usepackage{graphicx}
\graphicspath{{images/}}
\usepackage{subfig}
\usepackage{tikz}
\usetikzlibrary{shapes,arrows}
\usepackage{pgfplots}
\pgfplotsset{compat=newest}
\pgfplotsset{plot coordinates/math parser=false}
\newlength\figureheight
\newlength\figurewidth
\pgfkeys{/pgf/number format/.cd,
set decimal separator={,\!},
1000 sep={\,},
}
\usepackage{ifthen}
\usepackage{ifpdf}
\ifpdf
\usepackage[pdftex]{hyperref}
\else
\usepackage{hyperref}
\fi
\usepackage{color}
\hypersetup{%
colorlinks=true,
linkcolor=black,
citecolor=black,
urlcolor=black}

\renewcommand{\baselinestretch}{1.05}
\usepackage{fancyhdr}
\pagestyle{fancy}
\fancyfoot{}
\fancyhead[LE,RO]{\bfseries\thepage}
\fancyhead[RE]{\bfseries\nouppercase{\leftmark}}
\fancyhead[LO]{\bfseries\nouppercase{\rightmark}}
\setlength{\headheight}{15pt}

\let\headruleORIG\headrule
\renewcommand{\headrule}{\color{black} \headruleORIG}
\renewcommand{\headrulewidth}{1.0pt}
\usepackage{colortbl}
\arrayrulecolor{black}

\fancypagestyle{plain}{
  \fancyhead{}
  \fancyfoot[C]{\thepage}
  \renewcommand{\headrulewidth}{0pt}
}

\makeatletter
\def\@textbottom{\vskip \z@ \@plus 1pt}
\let\@texttop\relax
\makeatother

\makeatletter
\def\cleardoublepage{\clearpage\if@twoside \ifodd\c@page\else%
  \hbox{}%
  \thispagestyle{empty}%
  \newpage%
  \if@twocolumn\hbox{}\newpage\fi\fi\fi}
\makeatother

\usepackage{amsthm}
\usepackage{amssymb,amsmath,bbm}
\usepackage{array}
\usepackage{bm}
\usepackage{multirow}
\usepackage[footnote]{acronym}

\newcommand*{\SET}[1]  {\ensuremath{\mathbf{#1}}}
\newcommand*{\VEC}[1]  {\ensuremath{\boldsymbol{#1}}}
\newcommand*{\FAM}[1]  {\ensuremath{\boldsymbol{#1}}}
\newcommand*{\MAT}[1]  {\ensuremath{\boldsymbol{#1}}}
\newcommand*{\OP}[1]  {\ensuremath{\mathrm{#1}}}
\newcommand*{\NORM}[1]  {\ensuremath{\left\|#1\right\|}}
\newcommand*{\DPR}[2]  {\ensuremath{\left \langle #1,#2 \right \rangle}}
\newcommand*{\calbf}[1]  {\ensuremath{\boldsymbol{\mathcal{#1}}}}
\newcommand*{\shift}[1]  {\ensuremath{\boldsymbol{#1}}}

\newcommand{\eqdef}{\stackrel{\mathrm{def}}{=}}
\newcommand{\argmax}{\operatornamewithlimits{argmax}}
\newcommand{\argmin}{\operatornamewithlimits{argmin}}
\newcommand{\ud}{\, \mathrm{d}}
\newcommand{\vect}{\text{Vect}}
\newcommand{\sinc}{\ensuremath{\mathrm{sinc}}}
\newcommand{\esp}{\ensuremath{\mathbb{E}}}
\newcommand{\hilbert}{\ensuremath{\mathcal{H}}}
\newcommand{\fourier}{\ensuremath{\mathcal{F}}}
\newcommand{\sgn}{\text{sgn}}
\newcommand{\intTT}{\int_{-T}^{T}}
\newcommand{\intT}{\int_{-\frac{T}{2}}^{\frac{T}{2}}}
\newcommand{\intinf}{\int_{-\infty}^{+\infty}}
\newcommand{\Sh}{\ensuremath{\boldsymbol{S}}}
\newcommand{\C}{\SET{C}}
\newcommand{\R}{\SET{R}}
\newcommand{\Z}{\SET{Z}}
\newcommand{\N}{\SET{N}}
\newcommand{\K}{\SET{K}}
\newcommand{\reel}{\mathcal{R}}
\newcommand{\imag}{\mathcal{I}}
\newcommand{\cmnr}{c_{m,n}^\reel}
\newcommand{\cmni}{c_{m,n}^\imag}
\newcommand{\cnr}{c_{n}^\reel}
\newcommand{\cni}{c_{n}^\imag}
\newcommand{\tproto}{g}
\newcommand{\rproto}{\check{g}}
\newcommand{\LR}{\mathcal{L}_2(\SET{R})}
\newcommand{\LZ}{\ell_2(\SET{Z})}
\newcommand{\LZI}[1]{\ell_2(\SET{#1})}
\newcommand{\LZZ}{\ell_2(\SET{Z}^2)}
\newcommand{\diag}{\operatorname{diag}}
\newcommand{\noise}{z}
\newcommand{\Noise}{Z}
\newcommand{\filtnoise}{\zeta}
\newcommand{\tp}{g}
\newcommand{\rp}{\check{g}}
\newcommand{\TP}{G}
\newcommand{\RP}{\check{G}}
\newcommand{\dmin}{d_{\mathrm{min}}}
\newcommand{\Dmin}{D_{\mathrm{min}}}
\newcommand{\Image}{\ensuremath{\text{Im}}}
\newcommand{\Span}{\ensuremath{\text{Span}}}

\newtheoremstyle{break}
  {11pt}{11pt}%
  {\itshape}{}%
  {\bfseries}{}%
  {\newline}{}%
\theoremstyle{break}

%\theoremstyle{definition}
\newtheorem{definition}{Définition}[chapter]

%\theoremstyle{definition}
\newtheorem{theoreme}{Théorème}[chapter]

%\theoremstyle{remark}
\newtheorem{remarque}{Remarque}[chapter]

%\theoremstyle{plain}
\newtheorem{propriete}{Propriété}[chapter]
\newtheorem{exemple}{Exemple}[chapter]

\parskip=5pt
%\sloppy

\begin{document}

%%%%%%%%%%%%%%%%%%
%%% First page %%%
%%%%%%%%%%%%%%%%%%

\begin{titlepage}
\begin{center}

\includegraphics[width=0.6\textwidth]{logo-lille1-stdr}\\[1cm]

{\large Master Informatique spécialité TIIR}\\[0.5cm]

{\large Projet de Fin d'Etudes}\\[0.5cm]

% Title
\rule{\linewidth}{0.5mm} \\[0.4cm]
{ \huge \bfseries Déploiement d'une infrastructure pour enregistrer le trafic réseau de la salle TIIR \\[0.4cm] }
\rule{\linewidth}{0.5mm} \\[1.5cm]

% Author and supervisor
\noindent
\begin{minipage}{0.4\textwidth}
  \begin{flushleft} \large
    \emph{Auteurs :}\\
    M. Rémy \textsc{Debue}\\
    M. Kévin \textsc{Lanvin}\\
  \end{flushleft}
\end{minipage}%
\begin{minipage}{0.4\textwidth}
  \begin{flushright} \large
    \emph{Encadrants :} \\
    Julien \textsc{Iguchi-Catigny}
  \end{flushright}
\end{minipage}

\vfill

% Bottom of the page
{\large Version 0.1 du\\ \today}

\end{center}
\end{titlepage}

%%%%%%%%%%%%%%%%%%%%%%%%%%%%%
%%% Non-significant pages %%%
%%%%%%%%%%%%%%%%%%%%%%%%%%%%%

\frontmatter

\chapter*{Remerciements}
Nous souhaitons tout d'abord remercier Julien Iguchy-Cartigny, notre tuteur de projet qui nous a accompagné tout au long du PFE. Nos remerciements s'adressent également aux administrateurs du M5 qui nous ont acceuillis dans leur salle et dans la zone serveur et aidé à installer la nouvelle passerelle. Nous adressons nos remerciements également aux étudiants de la salle TIIR, pour leur patiente lors des coupures réseaux pendant que nous le configurions. Enfin, nous remercions les administrateurs du CRI qui se sont déplacés pour discuter de notre projet et de la nouvelle configuration du réseau.

\clearpage
\tableofcontents

\clearpage
\listoffigures

\clearpage
\chapter*{Liste des sigles et acronymes}
\begin{acronym}[CP-OFDMX] % Give the longest acronym here
\acro{CRI}{\emph{Centre de Ressources Informatique}}
\acro{DHCP}{\emph{Dynamic Host Configuration Protocol}}
\acro{DNS}{\emph{Domain Name Server}}
\acro{ELK}{\emph{Elasticsearch Logstash Kibana}}
\acro{PFE}{\emph{Projet de Fin d'Etudes}}
\acro{TIIR}{\emph{TIC, Infrastructures, Intégrité et Répartition}}
\end{acronym}

%%%%%%%%%%%%%%%%%%%%%%%%%%%%%%%%%%%%%%%%%%%%
%%% Content of the report and references %%%
%%%%%%%%%%%%%%%%%%%%%%%%%%%%%%%%%%%%%%%%%%%%

\mainmatter
\pagestyle{fancy}

\cleardoublepage

\chapter*{Introduction}
\addcontentsline{toc}{chapter}{Introduction}
\markboth{Introduction}{Introduction}
\label{chap:introduction}
%\minitoc

Le sujet "journalisation du trafic de la salle TIIR" trouve sa source dans le besoin d'autonomie des étudiants. En effet, Julien Iguchi-Cartigny, créateur de ce projet, souhaitait que l'on puisse s'affranchir du serveur proxy de l'université, afin de laisser plus de libertés pour la réalisation des travaux pratiques. Le proxy, serveur par lequel passait l'ensemble du trafic de la salle, permettait de filtrer les différentes connexions et ainsi empêcher les abus. Cependant, un tel filtrage ne peut pas être parfait, et certaines transmissions de données utiles dans le cadre des cours se trouvaient alors bloquées. Remplacer cette solution par une journalisation du trafic supprime ces contraintes pour les enseignants, mais permet de retrouver l'étudiant responsable en cas de comportement illicite (attaques sur le réseau d'entreprises, consultation de sites non-autorisés).

Ce projet de fin d'études se découpe alors en plusieurs parties. Tout d'abord, l'architecture du réseau de la salle TIIR a du être repensée. Il a fallu gérer l'installation des machines utiles pour notre projet, l'accès à internet de la salle et la répartition des adresses IP des postes de travail. Ensuite, un système a été mis en place pour enregistrer le trafic. La sauvegarde des données nécessitait un choix de technologie légère pour ne pas perturber le trafic. Les fichiers journaux sont envoyés vers une machine dédiée à leur indexation et leur analyse. Cette dernière proposera également une interface web permettant de visualiser les données plus rapidement.

%%% Local Variables:
%%% mode: latex
%%% TeX-master: "isae-report-template"
%%% End:

\chapter{Architecture}
\label{chap:Architecture}

\section{Analyse de l'existant}

A notre arrivée, le réseau de la salle TIIR était organisé autour d'un serveur proxy. Toutes les transmissions passaient à travers celui-ci, avant d'être chiffrées en utilisant IPSec. 
Ce système présentait plusieurs avantages. Le premier est la mise à disposition des administrateur d'un moyen de contrôler le trafic avec des règles de filtrage appliquées sur le proxy. Les activités illicites étaient ainsi rapidement détectées et bloquées automatiquement. Le second point positif est le fait que grâce à ce proxy, tout le trafic externe semblait provenir d'une seule et même machine, puisque l'adresse IP source de chaque paquet était l'adresse IP du proxy. L'inscription à de nouveaux services, comme l'accès à une bibliothèque digitale, en est ainsi facilité. Le chiffrement du trafic avec IPsec empêchait un quelconque vol des données sortant de l'université.

Néanmoins cette solution fût remise en question. Tout d'abord, l'absence de documentation sur cette architecture posait problème. Les technologies utilisées (le proxy et IPsec) furent laborieuses à mettre en place, et la maintenance s'annonçait ardue sans schéma ou fichiers de configuration à disposition. La notion de proxy posait également problème. En effet, le trafic était automatiquement filtré. Aucune confiance n'était accordée aux étudiants, et leurs marges de manoeuvre, lorsqu'ils souhaitaient sortir des sentiers battus était plus que limitée. Impossible de lancer la commande nmap qui permet entre autre de scanner les ports d'un poste, dans un tel réseau.

Une nouvelle solution a donc été envisagée : la journalisation du trafic. après la suppression du proxy, la liberté des étudiants augmenterait considérablement. Le trafic serait chiffré avant même d'arriver sur les postes des administrateurs de l'université, afin de les empêcher d'analyser les communications. Néanmoins la sécurité pénale de l'encadrant serait assurée grâce aux journaux que l'on pourrait consulter en cas de problème grave. De plus, cette structure pourrait servir de base pour des projets MOCAD. En effet le nombre important de communications enregistrées leur permettrait de s'entraîner à l'extraction d'informations dans de grands volumes de données.


\section{Conception de la nouvelle architecture}

\subsection{Accès à Internet}

Pour avoir accès à Internet depuis la salle TIIR, il faut mettre en place une passerelle. Cette passerelle est reliée d'un côté à la salle TIIR via un switch. Elle possède une adresse IP privée qui lui permet de communiquer avec tout le sous-réseau. De l'autre côté, cette passerelle possède une adresse publique. Cela signifie qu'elle peut communiquer sur Internet. Néanmoins, la topologie du réseau nous oblige à relier cette passerelle aux machines du CRI avant d'accéder à Internet (voir figure \ref{etape1}). 


\begin{figure}[!h]
\def\svgwidth{\columnwidth}
\input{images/etape1.pdf_tex}
\caption{Installation de la passerelle : Stargate}
\label{etape1}
\end{figure}

A ce stade, l'accès au réseau mondial est garanti, mais il reste deux problèmes. Le premier est la présence des machines du CRI sur la route des paquets sortants. Les administrateurs peuvent facilement lire, et bloquer les communications en provenance de la salle TIIR. Le second est que tout le trafic sortant a pour adresse IP la machine du CRI. Cela signifie que si un étudiant a un comportement illicite sur le réseau, l'université peut en être tenue pour responsable.

Pour résoudre ces problèmes, la solution retenue fût d'ajouter un VPS. Cela fixe une adresse IP source de toutes les communication sortantes différentes de celle de l'université. Celle-ci n'est donc plus directement l'émettrice des paquets interdits. Cela permet également de créer un tunnel sécurisé entre notre passerelle et le VPS (voir figure \ref{etape2}). Ainsi, les communications que le CRI retransmet sont chiffrées. Les administrateurs ne peuvent donc plus surveiller le trafic sortant de la salle TIIR, et encore moins le filtrer.

\begin{figure}[!h]
\centering
\def\svgwidth{\columnwidth}
\input{images/etape2.pdf_tex}
\caption{Mise en place d'un tunnel sécurisé}
\label{etape2}
\end{figure}

L'accès à Internet est maintenant garanti. Le CRI ne peut plus surveiller les communication et l'adresse de sortie n'est plus une adresse appartenant à l'université. Cette solution remplit donc les objectifs fixés auparavant.


\subsection{Architecture pour l'enregistrement des données}

Comme nous l'avons vu plus haut, le trafic de la salle TIIR dans son intégralité passe par la passerelle. C'est donc sur cette machine que l'enregistrement des données doit s'effectuer. Cependant, l'indexation et l'enregistrement des données consomme beaucoup de ressources. C'est pourquoi nous avons décidé d'utiliser une seconde machine pour effectuer cette tâche. Cette dernière fait partie du même sous-réseau que la salle TIIR. Ainsi, une fois les données analysées, un dashboard disponible sur une interface web permettra à tous les étudiants de constater certaines métriques sur les données transmises (voir figure \ref{etape3}).

\begin{figure}[!h]
\centering
\def\svgwidth{\columnwidth}
\input{images/etape3.pdf_tex}
\caption{Ajout d'une machine de traitement des journaux : Prism}
\label{etape3}
\end{figure}

Cette architecture est la dernière itération de notre raisonnement. Elle permet de contourner tous les problèmes inhérents à l'ancienne solution. Le CRI ne peut plus voir ni filtrer les paquets des étudiants, l'adresse de sortie n'est pas une adresse de l'université, l'enregistrement des données permet de se couvrir juridiquement en cas de problème, et enfin la séparation entre la machine qui enregistre les données et celle qui les traite permet un gain de performances non négligeable.


\section{Implémentation de la nouvelle architecture}

Nous allons présenter ici l'implémentation de l'architecture, qui permettra à tous les étudiants d'accéder à Internet, mais également à la machine qui traite les données d'être connectée à la passerelle.

\subsection{Accès à Internet}

La première étape pour garantir l'accès à Internet fût d'installer la passerelle Stargate. Nous avons utilisé un serveur disponible dans le bâtiment, sur lequel nous avons installé un Debian, un système d'exploitation Linux libre. Nous l'avons ensuite relié à la salle TIIR pour le configurer à distance. Le défi ici était d'installer cette machine en garantissant une interruption du service minimale. Nous avons donc conçu le plan d'adressage, le DHCP, le DNS et le tunnel VPN avant d'effectuer ce changement. Le VPS que nous avons loué appartient à OVH. Nous avons choisi ce prestataire car il offre une sécurité sur le matériel contre les attaques et garantit la disponibilité du service. C'est également l'un des partenaires de l'université.

\subsubsection{Le plan d'adressage}

La passerelle possède deux interfaces. La première est reliée à la salle TIIR, et la seconde est reliée indirectement au VPS. Pour la première, le choix de l'adressage a été assez simple. En effet, la salle TIIR appartenait déjà à un sous-réseau local : 10.1.1.0/24. La seconde est reliée à la passerelle du CRI. Une adresse IP lui a été distribuée : 134.206.225.12/32. Il était impératif dans notre configuration de faire passer le trafic via la passerelle du CRI qui avait pour adresse : 134.206.225.1/32. Cette dernière retransmettait alors les paquets vers le VPS OVH à l'adresse 137.74.40.238/32. La figure \ref{plan_adressage} résume ces adresses. Les modifications d'adresses sont faites dans le fichier /etc/network/interfaces qui permet de donner une adresse à une interface réseau, ainsi que définir sa passerelle par défaut.

\begin{figure}[!h]
\centering
\def\svgwidth{\columnwidth}
\input{images/plan_adressage.pdf_tex}
\caption{Plan d'adressage du réseau}
\label{plan_adressage}
\end{figure}

\subsubsection{Le serveur DHCP et le serveur DNS}

Avant de placer notre passerelle à la place de l'existante, certaines modifications restaient à effectuer pour garantir l'accès à Internet des utilisateurs de la salle TIIR. Tout d'abord, les adresses des postes de la salle ne sont pas fixes. Elles appartiennent toutes au sous-réseau 10.1.1.0/24 mais sont distribuées à la demande par un serveur DHCP qui se trouve sur la passerelle en service. Il faut également installer un serveur DNS qui permettra de résoudre les noms de domaines demandés par les étudiants, c'est à dire transformer les URL en adresses IP. S'il ne connait pas le domaine demandé, le serveur DNS devra rediriger la requête vers les serveurs DNS fournis par OVH. 

Pour implémenter cette solution, nous avons choisi d'utiliser dnsmasq. Cet outil permet en effet de regrouper les fonctions de DNS et de DHCP. La configuration choisie distribue des adresses entre 10.1.1.100 et 10.1.1.199. Cela permet de s'assurer une plage d'adresses inutilisées en cas d'ajout de machines qui nécessiteraient une adresse IP statique.

\subsubsection{Le tunnel VPN}

Afin d'éviter que le CRI ne filtre le trafic de la salle TIIR, celui-ci est chiffré par Stargate et déchiffré ensuite par le VPS OVH. Pour réaliser ce chiffrement, nous avons utilisé un tunnel VPN en point à point. Quand un étudiant émet un paquet depuis la salle TIIR, il arrive sur l'interface privée de la passerelle Stargate. Les règles de routage redirigent ensuite ce paquet vers une interface virtuelle (ici tun0). Cette interface est en fait reliée au programme OpenVPN qui chiffre le contenu du paquet. Puis, une seconde règle de routage redirige alors le paquet chiffré vers l'interface publique de la passerelle Stargate, qui expédie les données vers le CRI sur le port 1194, le port par  défaut pour les tunnels VPN. Nous avons demandé au CRI de retransmettre ces paquets vers le VPS OVH qui va effectuer le processus inverse. Les données seront redirigées vers une interface virtuelle reliée à OpenVPN, qui déchiffre le contenu des paquets et les renvoyer vers Internet. Les mêmes manipulations sont réalisées dans l'autre sens lorsque le serveur contacté par l'utilisateur répond. Les données sont chiffrées avec une clef secrète générée au préalable et transmise au VPS via un canal sécurisé. La figure \ref{tunnel} illustre ce principe.

\begin{figure}[!h]
\centering
\def\svgwidth{\columnwidth}
\input{images/tunnel.pdf_tex}
\caption{Fonctionnement du tunnel VPN}
\label{tunnel}
\end{figure}

\subsection{Enregistrement des données}

Afin d'enregistrer les données, nous avons ajouté une seconde machine au sous-réseau de la salle TIIR : Prism. Cette machine devra avoir une adresse IP statique pour être facilement accessible depuis le sous-réseau. Nous lui avons attribué l'adresse 10.1.1.2/32. Cette machine a été reformatée. Nous avons également installé un Debian. Cette machine a deux rôles. Elle doit indexer et enregistrer les données concernant le trafic de la salle TIIR mais également proposer une interface web pour faciliter la lecture de ces données. La figure \ref{architecture_finale} montre l'architecture finale du réseau.

\begin{figure}[!h]
\centering
\def\svgwidth{\columnwidth}
\input{images/architecture_finale.pdf_tex}
\caption{Architecture du réseau de la salle TIIR}
\label{architecture_finale}
\end{figure}


%%% Local Variables:
%%% mode: latex
%%% TeX-master: "isae-report-template"
%%% End:

\chapter{Notre projet et ses apports}
\label{sec:unchapitre}

Durant notre projet, plusieurs facettes du réseau ont été modifiées. En effet, l'architecture a été refaite, une nouvelle machine est apparue pour l'accès à Internet (par OVH). Nous avons également mis en place deux machines, une qui sert de passerelle, c'est à dire de faire le lien entre la salle TIIR et l'accès Internet à OVH. La deuxième permet d'analyser et de faire de la journalisation (ou log) du trafic sur le réseau.

\begin{figure}[htp!]
  \centering
  \setlength\figureheight{7cm}
  \setlength\figurewidth{9cm}
  % This file was created by matlab2tikz v0.2.2.
% Copyright (c) 2008--2012, Nico Schlömer <nico.schloemer@gmail.com>
% All rights reserved.
% 
% 
% 

% defining custom colors
\definecolor{mycolor1}{rgb}{0,0.75,0.75}

\begin{tikzpicture}

\begin{axis}[%
view={0}{90},
width=\figurewidth,
height=\figureheight,
scale only axis,
xmin=2, xmax=4.5,
xlabel={$\eta$},
xmajorgrids,
ymin=0.5, ymax=1,
ylabel={$d_{\text{min}}^2$},
ymajorgrids,
legend cell align=left,
legend style={align=left}]
\addplot [
color=black,
dashed,
mark=asterisk,
mark options={solid}
]
coordinates{
 (2,1)(2.1,1)(2.2,1)(2.3,1)(2.4,1)(2.5,1)(2.6,0.937749781479547)(2.7,0.890900393128398)(2.8,0.864988513955105)(2.9,0.827013168393703)(3,0.811347612650328)(3.1,0.792559278041243)(3.2,0.765840563467819)(3.3,0.749680961469385)(3.4,0.741947149227874)(3.5,0.740609493518419)(3.6,0.732128087463441)(3.7,0.717775843626632)(3.8,0.699687461812158)(3.9,0.685018622769455)(4,0.673439611642851)(4.1,0.664624248264608)(4.2,0.658255928882634)(4.3,0.641702335270489)(4.4,0.608326504614558)(4.5,0.580489221369454) 
};
\addlegendentry{$\alpha\text{ =  0\%}$};

\addplot [
color=black,
dashed,
mark=x,
mark options={solid}
]
coordinates{
 (2,1)(2.1,1)(2.2,1)(2.3,1)(2.4,0.958561324724996)(2.5,0.900812804739278)(2.6,0.859608621629443)(2.7,0.828484932127753)(2.8,0.812298837741994)(2.9,0.778916291864501)(3,0.758500630955482)(3.1,0.748375165853317)(3.2,0.745960208532468)(3.3,0.738441167434538)(3.4,0.715506361296671)(3.5,0.696927131434508)(3.6,0.682276848692725)(3.7,0.671128156410174)(3.8,0.663062783265717)(3.9,0.657680299791254)(4,0.621142740976429)(4.1,0.589786339121755)(4.2,0.564530571776849)(4.3,0.54483432747474)(4.4,0.53008799514765)(4.5,0.519641830384595) 
};
\addlegendentry{$\alpha\text{ = 10\%}$};

\addplot [
color=black,
dashed,
mark=triangle,
mark options={solid}
]
coordinates{
 (2,1)(2.1,1)(2.2,1)(2.3,0.966145915091813)(2.4,0.907589260275562)(2.5,0.862273165052718)(2.6,0.833762738286283)(2.7,0.797262289343802)(2.8,0.774689700869446)(2.9,0.763077871790574)(3,0.759584455148894)(3.1,0.735410358863577)(3.2,0.713220246811223)(3.3,0.695713299974315)(3.4,0.682371019886023)(3.5,0.672682085917092)(3.6,0.6661550402729)(3.7,0.644666127799479)(3.8,0.610083129739041)(3.9,0.582172698611821)(4,0.560333265725228)(4.1,0.543883933286703)(4.2,0.532098369213191)(4.3,0.524242326405)(4.4,0.519608701974017)(4.5,0.517545187250875) 
};
\addlegendentry{$\alpha\text{ = 20\%}$};

\addplot [
color=black,
dashed,
mark=triangle,
mark options={solid,,rotate=180}
]
coordinates{
 (2,1)(2.1,1)(2.2,0.995488894312993)(2.3,0.930050749246739)(2.4,0.882604857341179)(2.5,0.840148695151764)(2.6,0.807621264874927)(2.7,0.787889977099099)(2.8,0.777972678915356)(2.9,0.750463202108443)(3,0.726620292578349)(3.1,0.707917379352703)(3.2,0.693763185722015)(3.3,0.683575144048861)(3.4,0.676795290182409)(3.5,0.663350261880571)(3.6,0.627666127013326)(3.7,0.598755039468926)(3.8,0.575986310488554)(3.9,0.558651995817327)(4,0.546003746104731)(4.1,0.537291509323841)(4.2,0.531798375059385)(4.3,0.528867181690889)(4.4,0.527917002741411)(4.5,0.528450017604181) 
};
\addlegendentry{$\alpha\text{ = 30\%}$};

\addplot [
color=black,
dashed,
mark=o,
mark options={solid}
]
coordinates{
 (2,1)(2.1,1)(2.2,1)(2.3,1)(2.4,1)(2.5,0.995096871086856)(2.6,0.937749790013923)(2.7,0.890900391028178)(2.8,0.864988509535523)(2.9,0.827013167946275)(3,0.811347609462027)(3.1,0.79255927917077)(3.2,0.765840564829299)(3.3,0.749680963181722)(3.4,0.741947149533667)(3.5,0.740609492450166)(3.6,0.732128080624777)(3.7,0.71777584554089)(3.8,0.699687463368726)(3.9,0.681193180471954)(4,0.640212533267028)(4.1,0.617585040920557)(4.2,0.608519007405809)(4.3,0.608298095410932)(4.4,0.608326494076335)(4.5,0.580489212682311) 
};
\addlegendentry{Mazo};

\end{axis}
\end{tikzpicture}%
  \caption{Exemple de courbe TikZ.}
  \label{fig:courbe-tikz}
\end{figure}

\section{Nouvelle architecture du réseau}
Lorem ipsum dolor sit amet, consectetur adipiscing elit. Sed non risus. Suspendisse lectus tortor, dignissim sit amet, adipiscing nec, ultricies sed, dolor. Cras elementum ultrices diam. Maecenas ligula massa, varius a, semper congue, euismod non, mi. Proin porttitor, orci nec nonummy molestie, enim est eleifend mi, non fermentum diam nisl sit amet erat. Duis semper. Duis arcu massa, scelerisque vitae, consequat in, pretium a, enim. Pellentesque congue. Ut in risus volutpat libero pharetra tempor. Cras vestibulum bibendum augue. Praesent egestas leo in pede. Praesent blandit odio eu enim. Pellentesque sed dui ut augue blandit sodales. Vestibulum ante ipsum primis in faucibus orci luctus et ultrices posuere cubilia Curae; Aliquam nibh. Mauris ac mauris sed pede pellentesque fermentum. Maecenas adipiscing ante non diam sodales hendrerit. Ut velit mauris, egestas sed, gravida nec, ornare ut, mi. Aenean ut orci vel massa suscipit pulvinar. Nulla sollicitudin. Fusce varius, ligula non tempus aliquam, nunc turpis ullamcorper nibh, in tempus sapien eros vitae ligula. Pellentesque rhoncus nunc et augue. Integer id felis. Curabitur aliquet pellentesque diam. Integer quis metus vitae elit lobortis egestas. Lorem ipsum dolor sit amet, consectetuer adipiscing elit. Morbi vel erat non mauris convallis vehicula. Nulla et sapien. Integer tortor tellus, aliquam faucibus, convallis id, congue eu, quam. Mauris ullamcorper felis vitae erat. Proin feugiat, augue non elementum posuere, metus purus iaculis lectus, et tristique ligula justo vitae magna. Aliquam convallis sollicitudin purus. Praesent aliquam, enim at fermentum mollis, ligula massa adipiscing nisl, ac euismod nibh nisl eu lectus. Fusce vulputate sem at sapien. Vivamus leo. Aliquam euismod libero eu enim. Nulla nec felis sed leo placerat imperdiet. Aenean suscipit nulla in justo. Suspendisse cursus rutrum augue. Nulla tincidunt tincidunt mi. Curabitur iaculis, lorem vel rhoncus faucibus, felis magna fermentum augue, et ultricies lacus lorem varius purus. Curabitur eu amet.

\subsection{Machine à OVH et les machines à la fac}
Lorem ipsum dolor sit amet, consectetuer adipiscing elit. Morbi vel erat non mauris convallis vehicula. Nulla et sapien. Integer tortor tellus, aliquam faucibus, convallis id, congue eu, quam. Mauris ullamcorper felis vitae erat. Proin feugiat, augue non elementum posuere, metus purus iaculis lectus, et tristique ligula justo vitae magna. Aliquam convallis sollicitudin purus. Praesent aliquam, enim at fermentum mollis, ligula massa adipiscing nisl, ac euismod nibh nisl eu lectus. Fusce vulputate sem at sapien. Vivamus leo. Aliquam euismod libero eu enim. Nulla nec felis sed leo placerat imperdiet. Aenean suscipit nulla in justo. Suspendisse cursus rutrum augue. Nulla tincidunt tincidunt mi. Curabitur iaculis, lorem vel rhoncus faucibus, felis magna fermentum augue, et ultricies lacus lorem varius purus. Curabitur eu amet.

\subsection{Suppression du passage au CRI}
Lorem ipsum dolor sit amet, consectetuer adipiscing elit. Morbi vel erat non mauris convallis vehicula. Nulla et sapien. Integer tortor tellus, aliquam faucibus, convallis id, congue eu, quam. Mauris ullamcorper felis vitae erat. Proin feugiat, augue non elementum posuere, metus purus iaculis lectus, et tristique ligula justo vitae magna. Aliquam convallis sollicitudin purus. Praesent aliquam, enim at fermentum mollis, ligula massa adipiscing nisl, ac euismod nibh nisl eu lectus. Fusce vulputate sem at sapien. Vivamus leo. Aliquam euismod libero eu enim. Nulla nec felis sed leo placerat imperdiet. Aenean suscipit nulla in justo. Suspendisse cursus rutrum augue. Nulla tincidunt tincidunt mi. Curabitur iaculis, lorem vel rhoncus faucibus, felis magna fermentum augue, et ultricies lacus lorem varius purus. Curabitur eu amet.

\begin{align}
H_{m,n,p,q} &= \DPR{\rproto_{p,q}}{\OP{H} \tproto_{m,n}}\\
&= \iint\limits_{\SET{R}^2} S_{\OP{H}}(f,\tau) \DPR{\rproto_{p,q}}{\OP{U}_{f,\tau} \tproto_{m,n}} \ud f \ud \tau.
\end{align}

\section{Choix et pertincence de ces technologies}
Lorem ipsum dolor sit amet, consectetur adipiscing elit. Sed non risus. Suspendisse lectus tortor, dignissim sit amet, adipiscing nec, ultricies sed, dolor. Cras elementum ultrices diam. Maecenas ligula massa, varius a, semper congue, euismod non, mi. Proin porttitor, orci nec nonummy molestie, enim est eleifend mi, non fermentum diam nisl sit amet erat. Duis semper. Duis arcu massa, scelerisque vitae, consequat in, pretium a, enim. Pellentesque congue. Ut in risus volutpat libero pharetra tempor. Cras vestibulum bibendum augue. Praesent egestas leo in pede. Praesent blandit odio eu enim. Pellentesque sed dui ut augue blandit sodales. Vestibulum ante ipsum primis in faucibus orci luctus et ultrices posuere cubilia Curae; Aliquam nibh. Mauris ac mauris sed pede pellentesque fermentum. Maecenas adipiscing ante non diam sodales hendrerit. Ut velit mauris, egestas sed, gravida nec, ornare ut, mi. Aenean ut orci vel massa suscipit pulvinar. Nulla sollicitudin. Fusce varius, ligula non tempus aliquam, nunc turpis ullamcorper nibh, in tempus sapien eros vitae ligula. Pellentesque rhoncus nunc et augue. Integer id felis. Curabitur aliquet pellentesque diam. Integer quis metus vitae elit lobortis egestas. Lorem ipsum dolor sit amet, consectetuer adipiscing elit. Morbi vel erat non mauris convallis vehicula. Nulla et sapien. Integer tortor tellus, aliquam faucibus, convallis id, congue eu, quam. Mauris ullamcorper felis vitae erat. Proin feugiat, augue non elementum posuere, metus purus iaculis lectus, et tristique ligula justo vitae magna. Aliquam convallis sollicitudin purus. Praesent aliquam, enim at fermentum mollis, ligula massa adipiscing nisl, ac euismod nibh nisl eu lectus. Fusce vulputate sem at sapien. Vivamus leo. Aliquam euismod libero eu enim. Nulla nec felis sed leo placerat imperdiet. Aenean suscipit nulla in justo. Suspendisse cursus rutrum augue. Nulla tincidunt tincidunt mi. Curabitur iaculis, lorem vel rhoncus faucibus, felis magna fermentum augue, et ultricies lacus lorem varius purus. Curabitur eu amet.

%%% Local Variables: 
%%% mode: latex
%%% TeX-master: "isae-report-template"
%%% End: 

\chapter{Analyse du réseau}
\label{sec:thirdchapter}

Lorem ipsum dolor sit amet, consectetur adipiscing elit. Sed non risus. Suspendisse lectus tortor, dignissim sit amet, adipiscing nec, ultricies sed, dolor. Cras elementum ultrices diam. Maecenas ligula massa, varius a, semper congue, euismod non, mi. Proin porttitor, orci nec nonummy molestie, enim est eleifend mi, non fermentum diam nisl sit amet erat. Duis semper. Duis arcu massa, scelerisque vitae,  convallis sollicitudin purus. Praesent aliquam, enim at fermentum mollis, ligula massa adipiscing nisl, ac euismod nibh nisl eu lectus. Fusce vulputate sem at sapien. Vivamus leo. Aliquam euismod libero eu enim. Nulla nec felis sed leo placerat imperdiet. Aenean suscipit nulla in justo. Suspendisse cursus rutrum augue. Nulla tincidunt tincidunt mi. Curabitur iaculis, lorem vel rhoncus faucibus, felis magna fermentum augue, et ultricies lacus lorem varius purus. Curabitur eu amet. Encore une citation \cite{Cadambe2008}.

\begin{figure}[htp!]
  \centering
  \setlength\figureheight{7cm}
  \setlength\figurewidth{9cm}
  % This file was created by matlab2tikz v0.2.2.
% Copyright (c) 2008--2012, Nico Schlömer <nico.schloemer@gmail.com>
% All rights reserved.
% 
% 
% 

% defining custom colors
\definecolor{mycolor1}{rgb}{0,0.75,0.75}

\begin{tikzpicture}

\begin{axis}[%
view={0}{90},
width=\figurewidth,
height=\figureheight,
scale only axis,
xmin=2, xmax=4.5,
xlabel={$\eta$},
xmajorgrids,
ymin=0.5, ymax=1,
ylabel={$d_{\text{min}}^2$},
ymajorgrids,
legend cell align=left,
legend style={align=left}]
\addplot [
color=black,
dashed,
mark=asterisk,
mark options={solid}
]
coordinates{
 (2,1)(2.1,1)(2.2,1)(2.3,1)(2.4,1)(2.5,1)(2.6,0.937749781479547)(2.7,0.890900393128398)(2.8,0.864988513955105)(2.9,0.827013168393703)(3,0.811347612650328)(3.1,0.792559278041243)(3.2,0.765840563467819)(3.3,0.749680961469385)(3.4,0.741947149227874)(3.5,0.740609493518419)(3.6,0.732128087463441)(3.7,0.717775843626632)(3.8,0.699687461812158)(3.9,0.685018622769455)(4,0.673439611642851)(4.1,0.664624248264608)(4.2,0.658255928882634)(4.3,0.641702335270489)(4.4,0.608326504614558)(4.5,0.580489221369454) 
};
\addlegendentry{$\alpha\text{ =  0\%}$};

\addplot [
color=black,
dashed,
mark=x,
mark options={solid}
]
coordinates{
 (2,1)(2.1,1)(2.2,1)(2.3,1)(2.4,0.958561324724996)(2.5,0.900812804739278)(2.6,0.859608621629443)(2.7,0.828484932127753)(2.8,0.812298837741994)(2.9,0.778916291864501)(3,0.758500630955482)(3.1,0.748375165853317)(3.2,0.745960208532468)(3.3,0.738441167434538)(3.4,0.715506361296671)(3.5,0.696927131434508)(3.6,0.682276848692725)(3.7,0.671128156410174)(3.8,0.663062783265717)(3.9,0.657680299791254)(4,0.621142740976429)(4.1,0.589786339121755)(4.2,0.564530571776849)(4.3,0.54483432747474)(4.4,0.53008799514765)(4.5,0.519641830384595) 
};
\addlegendentry{$\alpha\text{ = 10\%}$};

\addplot [
color=black,
dashed,
mark=triangle,
mark options={solid}
]
coordinates{
 (2,1)(2.1,1)(2.2,1)(2.3,0.966145915091813)(2.4,0.907589260275562)(2.5,0.862273165052718)(2.6,0.833762738286283)(2.7,0.797262289343802)(2.8,0.774689700869446)(2.9,0.763077871790574)(3,0.759584455148894)(3.1,0.735410358863577)(3.2,0.713220246811223)(3.3,0.695713299974315)(3.4,0.682371019886023)(3.5,0.672682085917092)(3.6,0.6661550402729)(3.7,0.644666127799479)(3.8,0.610083129739041)(3.9,0.582172698611821)(4,0.560333265725228)(4.1,0.543883933286703)(4.2,0.532098369213191)(4.3,0.524242326405)(4.4,0.519608701974017)(4.5,0.517545187250875) 
};
\addlegendentry{$\alpha\text{ = 20\%}$};

\addplot [
color=black,
dashed,
mark=triangle,
mark options={solid,,rotate=180}
]
coordinates{
 (2,1)(2.1,1)(2.2,0.995488894312993)(2.3,0.930050749246739)(2.4,0.882604857341179)(2.5,0.840148695151764)(2.6,0.807621264874927)(2.7,0.787889977099099)(2.8,0.777972678915356)(2.9,0.750463202108443)(3,0.726620292578349)(3.1,0.707917379352703)(3.2,0.693763185722015)(3.3,0.683575144048861)(3.4,0.676795290182409)(3.5,0.663350261880571)(3.6,0.627666127013326)(3.7,0.598755039468926)(3.8,0.575986310488554)(3.9,0.558651995817327)(4,0.546003746104731)(4.1,0.537291509323841)(4.2,0.531798375059385)(4.3,0.528867181690889)(4.4,0.527917002741411)(4.5,0.528450017604181) 
};
\addlegendentry{$\alpha\text{ = 30\%}$};

\addplot [
color=black,
dashed,
mark=o,
mark options={solid}
]
coordinates{
 (2,1)(2.1,1)(2.2,1)(2.3,1)(2.4,1)(2.5,0.995096871086856)(2.6,0.937749790013923)(2.7,0.890900391028178)(2.8,0.864988509535523)(2.9,0.827013167946275)(3,0.811347609462027)(3.1,0.79255927917077)(3.2,0.765840564829299)(3.3,0.749680963181722)(3.4,0.741947149533667)(3.5,0.740609492450166)(3.6,0.732128080624777)(3.7,0.71777584554089)(3.8,0.699687463368726)(3.9,0.681193180471954)(4,0.640212533267028)(4.1,0.617585040920557)(4.2,0.608519007405809)(4.3,0.608298095410932)(4.4,0.608326494076335)(4.5,0.580489212682311) 
};
\addlegendentry{Mazo};

\end{axis}
\end{tikzpicture}%
  \caption{Exemple de courbe TikZ.}
  \label{fig:courbe-tikz}
\end{figure}

\section{Detection des paquets}
Lorem ipsum dolor sit amet, consectetur adipiscing elit. Sed non risus. Suspendisse lectus tortor, dignissim sit amet, adipiscing nec, ultricies sed, dolor. Cras elementum ultrices diam. Maecenas ligula massa, varius a, semper congue, euismod non, mi. Proin porttitor, orci nec nonummy molestie, enim est eleifend mi, non fermentum diam nisl sit amet erat. Duis semper. Duis arcu massa, scelerisque vitae, consequat in, pretium a, enim. Pellentesque congue. Ut in risus volutpat libero pharetra tempor. Cras vestibulum bibendum augue. Praesent egestas leo in pede. Praesent blandit odio eu enim. Pellentesque sed dui ut augue blandit sodales. Vestibulum ante ipsum primis in faucibus orci luctus et ultrices posuere cubilia Curae; Aliquam nibh. Mauris ac mauris sed pede pellentesque fermentum. Maecenas adipiscing ante non diam sodales hendrerit. Ut velit mauris, egestas sed, gravida nec, ornare ut, mi. Aenean ut orci vel massa suscipit pulvinar. Nulla sollicitudin. Fusce varius, ligula non tempus aliquam, nunc turpis ullamcorper nibh, in tempus sapien eros vitae ligula. Pellentesque rhoncus nunc et augue. Integer id felis. Curabitur aliquet pellentesque diam. Integer quis metus vitae elit lobortis egestas. Lorem ipsum dolor sit amet, consectetuer adipiscing elit. Morbi vel erat non mauris convallis vehicula. Nulla et sapien. Integer tortor tellus, aliquam faucibus, convallis id, congue eu, quam. Mauris ullamcorper felis vitae erat. Proin feugiat, augue non elementum posuere, metus purus iaculis lectus, et tristique ligula justo vitae magna. Aliquam convallis sollicitudin purus. Praesent aliquam, enim at fermentum mollis, ligula massa adipiscing nisl, ac euismod nibh nisl eu lectus. Fusce vulputate sem at sapien. Vivamus leo. Aliquam euismod libero eu enim. Nulla nec felis sed leo placerat imperdiet. Aenean suscipit nulla in justo. Suspendisse cursus rutrum augue. Nulla tincidunt tincidunt mi. Curabitur iaculis, lorem vel rhoncus faucibus, felis magna fermentum augue, et ultricies lacus lorem varius purus. Curabitur eu amet.

\subsection{A voir}
Lorem ipsum dolor sit amet, consectetuer adipiscing elit. Morbi vel erat non mauris convallis vehicula. Nulla et sapien. Integer tortor tellus, aliquam faucibus, convallis id, congue eu, quam. Mauris ullamcorper felis vitae erat. Proin feugiat, augue non elementum posuere, metus purus iaculis lectus, et tristique ligula justo vitae magna. Aliquam convallis sollicitudin purus. Praesent aliquam, enim at fermentum mollis, ligula massa adipiscing nisl, ac euismod nibh nisl eu lectus. Fusce vulputate sem at sapien. Vivamus leo. Aliquam euismod libero eu enim. Nulla nec felis sed leo placerat imperdiet. Aenean suscipit nulla in justo. Suspendisse cursus rutrum augue. Nulla tincidunt tincidunt mi. Curabitur iaculis, lorem vel rhoncus faucibus, felis magna fermentum augue, et ultricies lacus lorem varius purus. Curabitur eu amet.

\subsection{A voir 2}
Lorem ipsum dolor sit amet, consectetuer adipiscing elit. Morbi vel erat non mauris convallis vehicula. Nulla et sapien. Integer tortor tellus, aliquam faucibus, convallis id, congue eu, quam. Mauris ullamcorper felis vitae erat. Proin feugiat, augue non elementum posuere, metus purus iaculis lectus, et tristique ligula justo vitae magna. Aliquam convallis sollicitudin purus. Praesent aliquam, enim at fermentum mollis, ligula massa adipiscing nisl, ac euismod nibh nisl eu lectus. Fusce vulputate sem at sapien. Vivamus leo. Aliquam euismod libero eu enim. Nulla nec felis sed leo placerat imperdiet. Aenean suscipit nulla in justo. Suspendisse cursus rutrum augue. Nulla tincidunt tincidunt mi. Curabitur iaculis, lorem vel rhoncus faucibus, felis magna fermentum augue, et ultricies lacus lorem varius purus. Curabitur eu amet.

\begin{align}
H_{m,n,p,q} &= \DPR{\rproto_{p,q}}{\OP{H} \tproto_{m,n}}\\
&= \iint\limits_{\SET{R}^2} S_{\OP{H}}(f,\tau) \DPR{\rproto_{p,q}}{\OP{U}_{f,\tau} \tproto_{m,n}} \ud f \ud \tau.
\end{align}

\section{Log du trafic}
Lorem ipsum dolor sit amet, consectetur adipiscing elit. Sed non risus. Suspendisse lectus tortor, dignissim sit amet, adipiscing nec, ultricies sed, dolor. Cras elementum ultrices diam. Maecenas ligula massa, varius a, semper congue, euismod non, mi. Proin porttitor, orci nec nonummy molestie, enim est eleifend mi, non fermentum diam nisl sit amet erat. Duis semper. Duis arcu massa, scelerisque vitae, consequat in, pretium a, enim. Pellentesque congue. Ut in risus volutpat libero pharetra tempor. Cras vestibulum bibendum augue. Praesent egestas leo in pede. Praesent blandit odio eu enim. Pellentesque sed dui ut augue blandit sodales. Vestibulum ante ipsum primis in faucibus orci luctus et ultrices posuere cubilia Curae; Aliquam nibh. Mauris ac mauris sed pede pellentesque fermentum. Maecenas adipiscing ante non diam sodales hendrerit. Ut velit mauris, egestas sed, gravida nec, ornare ut, mi. Aenean ut orci vel massa suscipit pulvinar. Nulla sollicitudin. Fusce varius, ligula non tempus aliquam, nunc turpis ullamcorper nibh, in tempus sapien eros vitae ligula. Pellentesque rhoncus nunc et augue. Integer id felis. Curabitur aliquet pellentesque diam. Integer quis metus vitae elit lobortis egestas. Lorem ipsum dolor sit amet, consectetuer adipiscing elit. Morbi vel erat non mauris convallis vehicula. Nulla et sapien. Integer tortor tellus, aliquam faucibus, convallis id, congue eu, quam. Mauris ullamcorper felis vitae erat. Proin feugiat, augue non elementum posuere, metus purus iaculis lectus, et tristique ligula justo vitae magna. Aliquam convallis sollicitudin purus. Praesent aliquam, enim at fermentum mollis, ligula massa adipiscing nisl, ac euismod nibh nisl eu lectus. Fusce vulputate sem at sapien. Vivamus leo. Aliquam euismod libero eu enim. Nulla nec felis sed leo placerat imperdiet. Aenean suscipit nulla in justo. Suspendisse cursus rutrum augue. Nulla tincidunt tincidunt mi. Curabitur iaculis, lorem vel rhoncus faucibus, felis magna fermentum augue, et ultricies lacus lorem varius purus. Curabitur eu amet.

%%% Local Variables: 
%%% mode: latex
%%% TeX-master: "isae-report-template"
%%% End: 

\chapter{Resolution de la problématique et apports}
\label{sec:fourthchapter}

Lorem ipsum dolor sit amet, consectetur adipiscing elit. Sed non risus. Suspendisse lectus tortor, dignissim sit amet, adipiscing nec, ultricies sed, dolor. Cras elementum ultrices diam. Maecenas ligula massa, varius a, semper congue, euismod non, mi. Proin porttitor, orci nec nonummy molestie, enim est eleifend mi, non fermentum diam nisl sit amet erat. Duis semper. Duis arcu massa, scelerisque vitae,  convallis sollicitudin purus. Praesent aliquam, enim at fermentum mollis, ligula massa adipiscing nisl, ac euismod nibh nisl eu lectus. Fusce vulputate sem at sapien. Vivamus leo. Aliquam euismod libero eu enim. Nulla nec felis sed leo placerat imperdiet. Aenean suscipit nulla in justo. Suspendisse cursus rutrum augue. Nulla tincidunt tincidunt mi. Curabitur iaculis, lorem vel rhoncus faucibus, felis magna fermentum augue, et ultricies lacus lorem varius purus. Curabitur eu amet. Encore une citation \cite{Cadambe2008}.

\begin{figure}[htp!]
  \centering
  \setlength\figureheight{7cm}
  \setlength\figurewidth{9cm}
  % This file was created by matlab2tikz v0.2.2.
% Copyright (c) 2008--2012, Nico Schlömer <nico.schloemer@gmail.com>
% All rights reserved.
% 
% 
% 

% defining custom colors
\definecolor{mycolor1}{rgb}{0,0.75,0.75}

\begin{tikzpicture}

\begin{axis}[%
view={0}{90},
width=\figurewidth,
height=\figureheight,
scale only axis,
xmin=2, xmax=4.5,
xlabel={$\eta$},
xmajorgrids,
ymin=0.5, ymax=1,
ylabel={$d_{\text{min}}^2$},
ymajorgrids,
legend cell align=left,
legend style={align=left}]
\addplot [
color=black,
dashed,
mark=asterisk,
mark options={solid}
]
coordinates{
 (2,1)(2.1,1)(2.2,1)(2.3,1)(2.4,1)(2.5,1)(2.6,0.937749781479547)(2.7,0.890900393128398)(2.8,0.864988513955105)(2.9,0.827013168393703)(3,0.811347612650328)(3.1,0.792559278041243)(3.2,0.765840563467819)(3.3,0.749680961469385)(3.4,0.741947149227874)(3.5,0.740609493518419)(3.6,0.732128087463441)(3.7,0.717775843626632)(3.8,0.699687461812158)(3.9,0.685018622769455)(4,0.673439611642851)(4.1,0.664624248264608)(4.2,0.658255928882634)(4.3,0.641702335270489)(4.4,0.608326504614558)(4.5,0.580489221369454) 
};
\addlegendentry{$\alpha\text{ =  0\%}$};

\addplot [
color=black,
dashed,
mark=x,
mark options={solid}
]
coordinates{
 (2,1)(2.1,1)(2.2,1)(2.3,1)(2.4,0.958561324724996)(2.5,0.900812804739278)(2.6,0.859608621629443)(2.7,0.828484932127753)(2.8,0.812298837741994)(2.9,0.778916291864501)(3,0.758500630955482)(3.1,0.748375165853317)(3.2,0.745960208532468)(3.3,0.738441167434538)(3.4,0.715506361296671)(3.5,0.696927131434508)(3.6,0.682276848692725)(3.7,0.671128156410174)(3.8,0.663062783265717)(3.9,0.657680299791254)(4,0.621142740976429)(4.1,0.589786339121755)(4.2,0.564530571776849)(4.3,0.54483432747474)(4.4,0.53008799514765)(4.5,0.519641830384595) 
};
\addlegendentry{$\alpha\text{ = 10\%}$};

\addplot [
color=black,
dashed,
mark=triangle,
mark options={solid}
]
coordinates{
 (2,1)(2.1,1)(2.2,1)(2.3,0.966145915091813)(2.4,0.907589260275562)(2.5,0.862273165052718)(2.6,0.833762738286283)(2.7,0.797262289343802)(2.8,0.774689700869446)(2.9,0.763077871790574)(3,0.759584455148894)(3.1,0.735410358863577)(3.2,0.713220246811223)(3.3,0.695713299974315)(3.4,0.682371019886023)(3.5,0.672682085917092)(3.6,0.6661550402729)(3.7,0.644666127799479)(3.8,0.610083129739041)(3.9,0.582172698611821)(4,0.560333265725228)(4.1,0.543883933286703)(4.2,0.532098369213191)(4.3,0.524242326405)(4.4,0.519608701974017)(4.5,0.517545187250875) 
};
\addlegendentry{$\alpha\text{ = 20\%}$};

\addplot [
color=black,
dashed,
mark=triangle,
mark options={solid,,rotate=180}
]
coordinates{
 (2,1)(2.1,1)(2.2,0.995488894312993)(2.3,0.930050749246739)(2.4,0.882604857341179)(2.5,0.840148695151764)(2.6,0.807621264874927)(2.7,0.787889977099099)(2.8,0.777972678915356)(2.9,0.750463202108443)(3,0.726620292578349)(3.1,0.707917379352703)(3.2,0.693763185722015)(3.3,0.683575144048861)(3.4,0.676795290182409)(3.5,0.663350261880571)(3.6,0.627666127013326)(3.7,0.598755039468926)(3.8,0.575986310488554)(3.9,0.558651995817327)(4,0.546003746104731)(4.1,0.537291509323841)(4.2,0.531798375059385)(4.3,0.528867181690889)(4.4,0.527917002741411)(4.5,0.528450017604181) 
};
\addlegendentry{$\alpha\text{ = 30\%}$};

\addplot [
color=black,
dashed,
mark=o,
mark options={solid}
]
coordinates{
 (2,1)(2.1,1)(2.2,1)(2.3,1)(2.4,1)(2.5,0.995096871086856)(2.6,0.937749790013923)(2.7,0.890900391028178)(2.8,0.864988509535523)(2.9,0.827013167946275)(3,0.811347609462027)(3.1,0.79255927917077)(3.2,0.765840564829299)(3.3,0.749680963181722)(3.4,0.741947149533667)(3.5,0.740609492450166)(3.6,0.732128080624777)(3.7,0.71777584554089)(3.8,0.699687463368726)(3.9,0.681193180471954)(4,0.640212533267028)(4.1,0.617585040920557)(4.2,0.608519007405809)(4.3,0.608298095410932)(4.4,0.608326494076335)(4.5,0.580489212682311) 
};
\addlegendentry{Mazo};

\end{axis}
\end{tikzpicture}%
  \caption{Exemple de courbe TikZ.}
  \label{fig:courbe-tikz}
\end{figure}

\section{Les problèmes juridiques}
Lorem ipsum dolor sit amet, consectetur adipiscing elit. Sed non risus. Suspendisse lectus tortor, dignissim sit amet, adipiscing nec, ultricies sed, dolor. Cras elementum ultrices diam. Maecenas ligula massa, varius a, semper congue, euismod non, mi. Proin porttitor, orci nec nonummy molestie, enim est eleifend mi, non fermentum diam nisl sit amet erat. Duis semper. Duis arcu massa, scelerisque vitae, consequat in, pretium a, enim. Pellentesque congue. Ut in risus volutpat libero pharetra tempor. Cras vestibulum bibendum augue. Praesent egestas leo in pede. Praesent blandit odio eu enim. Pellentesque sed dui ut augue blandit sodales. Vestibulum ante ipsum primis in faucibus orci luctus et ultrices posuere cubilia Curae; Aliquam nibh. Mauris ac mauris sed pede pellentesque fermentum. Maecenas adipiscing ante non diam sodales hendrerit. Ut velit mauris, egestas sed, gravida nec, ornare ut, mi. Aenean ut orci vel massa suscipit pulvinar. Nulla sollicitudin. Fusce varius, ligula non tempus aliquam, nunc turpis ullamcorper nibh, in tempus sapien eros vitae ligula. Pellentesque rhoncus nunc et augue. Integer id felis. Curabitur aliquet pellentesque diam. Integer quis metus vitae elit lobortis egestas. Lorem ipsum dolor sit amet, consectetuer adipiscing elit. Morbi vel erat non mauris convallis vehicula. Nulla et sapien. Integer tortor tellus, aliquam faucibus, convallis id, congue eu, quam. Mauris ullamcorper felis vitae erat. Proin feugiat, augue non elementum posuere, metus purus iaculis lectus, et tristique ligula justo vitae magna. Aliquam convallis sollicitudin purus. Praesent aliquam, enim at fermentum mollis, ligula massa adipiscing nisl, ac euismod nibh nisl eu lectus. Fusce vulputate sem at sapien. Vivamus leo. Aliquam euismod libero eu enim. Nulla nec felis sed leo placerat imperdiet. Aenean suscipit nulla in justo. Suspendisse cursus rutrum augue. Nulla tincidunt tincidunt mi. Curabitur iaculis, lorem vel rhoncus faucibus, felis magna fermentum augue, et ultricies lacus lorem varius purus. Curabitur eu amet.

\subsection{A voir}
Lorem ipsum dolor sit amet, consectetuer adipiscing elit. Morbi vel erat non mauris convallis vehicula. Nulla et sapien. Integer tortor tellus, aliquam faucibus, convallis id, congue eu, quam. Mauris ullamcorper felis vitae erat. Proin feugiat, augue non elementum posuere, metus purus iaculis lectus, et tristique ligula justo vitae magna. Aliquam convallis sollicitudin purus. Praesent aliquam, enim at fermentum mollis, ligula massa adipiscing nisl, ac euismod nibh nisl eu lectus. Fusce vulputate sem at sapien. Vivamus leo. Aliquam euismod libero eu enim. Nulla nec felis sed leo placerat imperdiet. Aenean suscipit nulla in justo. Suspendisse cursus rutrum augue. Nulla tincidunt tincidunt mi. Curabitur iaculis, lorem vel rhoncus faucibus, felis magna fermentum augue, et ultricies lacus lorem varius purus. Curabitur eu amet.

\subsection{A voir 2}
Lorem ipsum dolor sit amet, consectetuer adipiscing elit. Morbi vel erat non mauris convallis vehicula. Nulla et sapien. Integer tortor tellus, aliquam faucibus, convallis id, congue eu, quam. Mauris ullamcorper felis vitae erat. Proin feugiat, augue non elementum posuere, metus purus iaculis lectus, et tristique ligula justo vitae magna. Aliquam convallis sollicitudin purus. Praesent aliquam, enim at fermentum mollis, ligula massa adipiscing nisl, ac euismod nibh nisl eu lectus. Fusce vulputate sem at sapien. Vivamus leo. Aliquam euismod libero eu enim. Nulla nec felis sed leo placerat imperdiet. Aenean suscipit nulla in justo. Suspendisse cursus rutrum augue. Nulla tincidunt tincidunt mi. Curabitur iaculis, lorem vel rhoncus faucibus, felis magna fermentum augue, et ultricies lacus lorem varius purus. Curabitur eu amet.

\begin{align}
H_{m,n,p,q} &= \DPR{\rproto_{p,q}}{\OP{H} \tproto_{m,n}}\\
&= \iint\limits_{\SET{R}^2} S_{\OP{H}}(f,\tau) \DPR{\rproto_{p,q}}{\OP{U}_{f,\tau} \tproto_{m,n}} \ud f \ud \tau.
\end{align}

\section{Indépendance des utilisateurs sur le réseau}
Lorem ipsum dolor sit amet, consectetur adipiscing elit. Sed non risus. Suspendisse lectus tortor, dignissim sit amet, adipiscing nec, ultricies sed, dolor. Cras elementum ultrices diam. Maecenas ligula massa, varius a, semper congue, euismod non, mi. Proin porttitor, orci nec nonummy molestie, enim est eleifend mi, non fermentum diam nisl sit amet erat. Duis semper. Duis arcu massa, scelerisque vitae, consequat in, pretium a, enim. Pellentesque congue. Ut in risus volutpat libero pharetra tempor. Cras vestibulum bibendum augue. Praesent egestas leo in pede. Praesent blandit odio eu enim. Pellentesque sed dui ut augue blandit sodales. Vestibulum ante ipsum primis in faucibus orci luctus et ultrices posuere cubilia Curae; Aliquam nibh. Mauris ac mauris sed pede pellentesque fermentum. Maecenas adipiscing ante non diam sodales hendrerit. Ut velit mauris, egestas sed, gravida nec, ornare ut, mi. Aenean ut orci vel massa suscipit pulvinar. Nulla sollicitudin. Fusce varius, ligula non tempus aliquam, nunc turpis ullamcorper nibh, in tempus sapien eros vitae ligula. Pellentesque rhoncus nunc et augue. Integer id felis. Curabitur aliquet pellentesque diam. Integer quis metus vitae elit lobortis egestas. Lorem ipsum dolor sit amet, consectetuer adipiscing elit. Morbi vel erat non mauris convallis vehicula. Nulla et sapien. Integer tortor tellus, aliquam faucibus, convallis id, congue eu, quam. Mauris ullamcorper felis vitae erat. Proin feugiat, augue non elementum posuere, metus purus iaculis lectus, et tristique ligula justo vitae magna. Aliquam convallis sollicitudin purus. Praesent aliquam, enim at fermentum mollis, ligula massa adipiscing nisl, ac euismod nibh nisl eu lectus. Fusce vulputate sem at sapien. Vivamus leo. Aliquam euismod libero eu enim. Nulla nec felis sed leo placerat imperdiet. Aenean suscipit nulla in justo. Suspendisse cursus rutrum augue. Nulla tincidunt tincidunt mi. Curabitur iaculis, lorem vel rhoncus faucibus, felis magna fermentum augue, et ultricies lacus lorem varius purus. Curabitur eu amet.
\section{Pédagogie concernant la sécurité réseau}
Lorem ipsum dolor sit amet, consectetur adipiscing elit. Sed non risus. Suspendisse lectus tortor, dignissim sit amet, adipiscing nec, ultricies sed, dolor. Cras elementum ultrices diam. Maecenas ligula massa, varius a, semper congue, euismod non, mi. Proin porttitor, orci nec nonummy molestie, enim est eleifend mi, non fermentum diam nisl sit amet erat. Duis semper. Duis arcu massa, scelerisque vitae, consequat in, pretium a, enim. Pellentesque congue. Ut in risus volutpat libero pharetra tempor. Cras vestibulum bibendum augue. Praesent egestas leo in pede. Praesent blandit odio eu enim. Pellentesque sed dui ut augue blandit sodales. Vestibulum ante ipsum primis in faucibus orci luctus et ultrices posuere cubilia Curae; Aliquam nibh. Mauris ac mauris sed pede pellentesque fermentum. Maecenas adipiscing ante non diam sodales hendrerit. Ut velit mauris, egestas sed, gravida nec, ornare ut, mi. Aenean ut orci vel massa suscipit pulvinar. Nulla sollicitudin. Fusce varius, ligula non tempus aliquam, nunc turpis ullamcorper nibh, in tempus sapien eros vitae ligula. Pellentesque rhoncus nunc et augue. Integer id felis. Curabitur aliquet pellentesque diam. Integer quis metus vitae elit lobortis egestas. Lorem ipsum dolor sit amet, consectetuer adipiscing elit. Morbi vel erat non mauris convallis vehicula. Nulla et sapien. Integer tortor tellus, aliquam faucibus, convallis id, congue eu, quam. Mauris ullamcorper felis vitae erat. Proin feugiat, augue non elementum posuere, metus purus iaculis lectus, et tristique ligula justo vitae magna. Aliquam convallis sollicitudin purus. Praesent aliquam, enim at fermentum mollis, ligula massa adipiscing nisl, ac euismod nibh nisl eu lectus. Fusce vulputate sem at sapien. Vivamus leo. Aliquam euismod libero eu enim. Nulla nec felis sed leo placerat imperdiet. Aenean suscipit nulla in justo. Suspendisse cursus rutrum augue. Nulla tincidunt tincidunt mi. Curabitur iaculis, lorem vel rhoncus faucibus, felis magna fermentum augue, et ultricies lacus lorem varius purus. Curabitur eu amet.

%%% Local Variables: 
%%% mode: latex
%%% TeX-master: "isae-report-template"
%%% End: 

\chapter*{Conclusion et perspectives}
\addcontentsline{toc}{chapter}{Conclusion}
\markboth{Conclusion}{Conclusion}
\label{sec:conclusion}

    Lorem ipsum dolor sit amet, consectetur adipiscing elit. Sed non risus. Suspendisse lectus tortor, dignissim sit amet, adipiscing nec, ultricies sed, dolor. Cras elementum ultrices diam. Maecenas ligula massa, varius a, semper congue, euismod non, mi. Proin porttitor, orci nec nonummy molestie, enim est eleifend mi, non fermentum diam nisl sit amet erat. Duis semper. Duis arcu massa, scelerisque vitae, consequat in, pretium a, enim. Pellentesque congue. Ut in risus volutpat libero pharetra tempor. Cras vestibulum bibendum augue. Praesent egestas leo in pede. Praesent blandit odio eu enim. Pellentesque sed dui ut augue blandit sodales. Vestibulum ante ipsum primis in faucibus orci luctus et ultrices posuere cubilia Curae; Aliquam nibh. Mauris ac mauris sed pede pellentesque fermentum. Maecenas adipiscing ante non diam sodales hendrerit. Ut velit mauris, egestas sed, gravida nec, ornare ut, mi. Aenean ut orci vel massa suscipit pulvinar. Nulla sollicitudin. Fusce varius, ligula non tempus aliquam, nunc turpis ullamcorper nibh, in tempus sapien eros vitae ligula. Pellentesque rhoncus nunc et augue. Integer id felis. Curabitur aliquet pellentesque diam. Integer quis metus vitae elit lobortis egestas. Lorem ipsum dolor sit amet, consectetuer adipiscing elit. Morbi vel erat non mauris convallis vehicula. Nulla et sapien. Integer tortor tellus, aliquam faucibus, convallis id, congue eu, quam. Mauris ullamcorper felis vitae erat. Proin feugiat, augue non elementum posuere, metus purus iaculis lectus, et tristique ligula justo vitae magna. Aliquam convallis sollicitudin purus. Praesent aliquam, enim at fermentum mollis, ligula massa adipiscing nisl, ac euismod nibh nisl eu lectus. Fusce vulputate sem at sapien. Vivamus leo. Aliquam euismod libero eu enim. Nulla nec felis sed leo placerat imperdiet. Aenean suscipit nulla in justo. Suspendisse cursus rutrum augue. Nulla tincidunt tincidunt mi. Curabitur iaculis, lorem vel rhoncus faucibus, felis magna fermentum augue, et ultricies lacus lorem varius purus. Curabitur eu amet.

%%% Local Variables: 
%%% mode: latex
%%% TeX-master: "isae-report-template"
%%% End: 



\appendix

\bibliographystyle{authoryear-fr}
\bibliography{references}

\clearpage

%%%%%%%%%%%%%%%%
%%% Abstract %%%
%%%%%%%%%%%%%%%%

\thispagestyle{empty}

\vspace*{\fill}
\noindent\rule[2pt]{\textwidth}{0.5pt}\\
{\textbf{Résumé ---}}
Lorem ipsum dolor sit amet, consectetur adipiscing elit. Sed non risus. Suspendisse lectus tortor, dignissim sit amet, adipiscing nec, ultricies sed, dolor. Cras elementum ultrices diam. Maecenas ligula massa, varius a, semper congue, euismod non, mi. Proin porttitor, orci nec nonummy molestie, enim est eleifend mi, non fermentum diam nisl sit amet erat. Duis semper. Duis arcu massa, scelerisque vitae, consequat in, pretium a, enim. Pellentesque congue. Ut in risus volutpat libero pharetra tempor. Cras vestibulum bibendum augue. Praesent egestas leo in pede. Praesent blandit odio eu enim. Pellentesque sed dui ut augue blandit sodales. Vestibulum ante ipsum primis in faucibus orci luctus et ultrices posuere cubilia Curae; Aliquam nibh. Mauris ac mauris sed pede pellentesque fermentum. Maecenas adipiscing ante non diam sodales hendrerit. Ut velit mauris, egestas sed, gravida nec, ornare ut, mi. Aenean ut orci vel massa suscipit pulvinar. Nulla sollicitudin. Fusce varius, ligula non tempus aliquam, nunc turpis ullamcorper nibh, in tempus sapien eros vitae ligula. Pellentesque rhoncus nunc et augue. Integer id felis.

{\textbf{Mots clés :}}
Lorem ipsum dolor sit amet, consectetur adipiscing elit. Sed non risus. Suspendisse lectus tortor.
\\
\noindent\rule[2pt]{\textwidth}{0.5pt}
\begin{center}
  Université de Lille 1\\
  Cité Scientifique\\
  59650 Villeneuve d'Ascq France
\end{center}
\vspace*{\fill}

\end{document}
